% !TeX root = ../main.tex
\clearpage
\section{Arsitektur Sistem}
\label{sec:architecture}
Gambar~\ref{fig:architecture} memperlihatkan container, protokol, dan kepemilikan storage sistem.
\architecturefigure

\subsection{Alur kerja sistem}
Alur sistem pada Gambar~\ref{fig:architecture} terdiri atas penerimaan data, pelayanan permintaan client, dan distribusi event. Ketiganya menggunakan data kanonik yang sama, tetapi berjalan secara terpisah agar keterlambatan sumber atau consumer tidak langsung menahan permintaan baca.

Pada penerimaan data, worker Aggregator melakukan polling ke BMKG dan PVMBG dengan kredensial masing-masing. Record divalidasi, dipetakan menjadi HazardEvent, dan dikorelasikan dengan warning tsunami bila relevan. Data baru atau perubahan konten disimpan bersama snapshot outbox dalam transaksi PostgreSQL. Record yang isinya tetap tidak menghasilkan event baru. Checkpoint maju setelah seluruh respons selesai ditangani, sehingga kegagalan di tengah proses masih dapat dipulihkan melalui pembacaan ulang. Rincian pemetaan dan transaksi dibahas pada Bagian~\ref{sec:p1}.

Pada pelayanan client, identitas ditukar menjadi token melalui Auth Service. Client kemudian mengirim permintaan ber-JWT ke Client API, yang memeriksa hak akses dan membatasi beban sebelum meneruskan query ke API internal Aggregator. Aggregator membaca Canonical Store tanpa memanggil sumber kembali. Client API memproyeksikan hasil sesuai scope dan menyertakan status kebaruan sumber. Pemisahan ini mempertahankan akses ke data tersimpan saat sumber bermasalah, dengan konsekuensi data dapat bersifat basi. Mekanismenya dijelaskan pada Bagian~\ref{sec:p2} dan~\ref{sec:p3}.

Pada distribusi event, relay membaca outbox yang sudah committed dan memublikasikan snapshot ke Kafka. Setelah ACK broker diterima, baris outbox ditandai published. Dashboard Updater, Notifier, dan Pemda Portal memakai consumer group berbeda serta penyimpanan SQLite masing-masing. Dengan demikian, setiap consumer dapat memproses stream dan mengejar ketertinggalan secara independen tanpa meminta ulang data ke Aggregator. Batas pengiriman ulang dan idempotensi dibahas pada Bagian~\ref{sec:p5}.

\subsection{Komponen dan kepemilikan data}
Sistem memiliki delapan service aplikasi apabila pemda-portal diaktifkan. PostgreSQL, Redis, dan Kafka merupakan container infrastruktur. Konfigurasi dasar menjalankan sepuluh container jangka panjang, ditambah satu \texttt{kafka-init} yang selesai setelah membuat topic; profile demo menambah pemda-portal. CLI dan load generator adalah client pengujian, bukan service penyimpan data.

\begin{longtable}{P{0.22\linewidth}P{0.43\linewidth}P{0.25\linewidth}}
\toprule Komponen & Tanggung jawab & Storage dan port\\\midrule\endhead
\texttt{bmkg-mock} & Menyajikan gempa dan warning terpisah; seed dan generator; memvalidasi API key BMKG. & Memori sendiri; 8081.\\
\texttt{pvmbg-mock} & Laporan vulkanik, delay, perubahan skema, dan outage yang dapat dipicu operator. & Memori sendiri; 8082.\\
\texttt{aggregator} & Polling, tolerant reader, pemetaan, korelasi, transaksi, query internal, dan relay outbox. & Pemilik PostgreSQL; 9000 internal.\\
\texttt{client-api} & Verifikasi JWT, scope, proyeksi, filter, pagination, dan proteksi beban. & Tanpa store persisten; 8080.\\
\texttt{auth-service} & Penerbit access token dan refresh token dan pengelola rotasi sesi. & Pemilik Redis; 8090.\\
\texttt{dashboard-updater} & View event terbaru per hazard, hanya menerima versi yang lebih baru. & SQLite sendiri; 8091.\\
\texttt{notifier} & Simulasi pengiriman alert SIAGA atau AWAS; dedup persisten. & SQLite sendiri; 8092.\\
\texttt{pemda-portal} & Subscriber tambahan untuk menunjukkan perluasan stream tanpa mengubah producer. & SQLite sendiri; 8093.\\
\bottomrule
\end{longtable}

Port host aplikasi dipublikasikan pada loopback \texttt{127.0.0.1}. Port database, Redis, Kafka, dan API internal Aggregator tidak dipublikasikan ke host pada konfigurasi normal. Container memakai named volume untuk store persisten; restart biasa tidak menghapus state.

\subsection{Protokol dan batas kepercayaan}
Aggregator mengambil data JSON melalui HTTP dengan kredensial yang berbeda untuk setiap instansi. BMKG memerlukan \texttt{X-BMKG-Key}, sedangkan PVMBG memerlukan Bearer token miliknya sendiri. Client menukar kredensial ke \texttt{/oauth/token}, kemudian mengirim JWT ke client-api. Client-api membaca lewat HTTP internal Aggregator dengan \texttt{X-Internal-Key}; ia tidak mempunyai kredensial PostgreSQL. Auth-service mengakses Redis melalui RESP dan password yang berbeda.

Relay mengirim envelope melalui protokol Kafka. Consumer mengambil stream dengan group masing-masing dan melakukan efek lokal di SQLite. Tidak ada pemanggilan fungsi maupun shared business package lintas service. Kontrak yang sama direpresentasikan secara lokal di setiap module. Komponen dalam Aggregator tetap boleh berkomunikasi lewat port berupa interface Go karena berada pada service pemilik yang sama.

Network \texttt{store\_net} menghubungkan Aggregator dan PostgreSQL; \texttt{edge\_net} menghubungkan client-api dan Aggregator; \texttt{source\_net} menghubungkan sumber dan poller; \texttt{auth\_net} menghubungkan auth-service dan Redis; \texttt{bus\_net} menghubungkan producer, broker, dan consumer. Profile test mempunyai akses yang diperlukan sebagai perangkat uji administratif, bukan jalur akses aplikasi produksi.
\coderef{docker-compose.yml}
